% Part: computability
% Chapter: recursive-functions
% Section: trees

\documentclass[../../../include/open-logic-section]{subfiles}

\begin{document}

% SEGMENT OLP-0221-S01

\olfileid{cmp}{rec}{tre}
\olsection{மரவுருக்கள்}

தொடர்களை எண்களாலும், அவற்றின் பண்புகளையும் அவற்றின்மீதான செயல்களையும்
அவற்றின் குறியீடுகள்மீதான முதன்மை மீள்வரையறை உறவுகள் மற்றும்
சார்புகளாலும் குறிக்க முடிவதைப் போலவே, சிலவேளைகளில் மரவுருக்களை
இயல் எண்களாகக் குறிப்பது பயனுள்ளது. இதைச் செய்ய தொடர்களையும் அவற்றின்
குறியீடுகளையும் பயன்படுத்துவோம். ஒரு மரவுரு, ஒரு தனிக் கணுவாக
(குறிச்சீட்டுடன் இருக்கலாம்) அல்லது பல உட்மரவுருக்களுடன் இணைக்கப்பட்ட
ஒரு கணுவாக (குறிச்சீட்டுடன் இருக்கலாம்) இருக்கலாம். அந்தக் கணு
மரவுருவின் \emph{வேர்} என்றும், அதனுடன் இணைக்கப்பட்ட உட்மரவுருக்கள்
அதன் \emph{உடனடி உட்மரவுருக்கள்} என்றும் அழைக்கப்படுகின்றன.

மரவுருக்களை $\tuple{k, d_1, \dots, d_k}$ என்ற தொடராக மீள்வரையறை
முறையில் குறியீடாக்குகிறோம்; இங்கே $k$ என்பது உடனடி
உட்மரவுருக்களின் எண்ணிக்கை; $d_1$, \dots,~$d_k$ என்பவை உடனடி
உட்மரவுருக்களின் குறியீடுகள். கணுக்கள் குறிச்சீட்டுகளைக் கொண்டிருந்தால்,
உடனடி உட்மரவுருக்களுக்குப் பின்னர் அவற்றைச் சேர்க்கலாம். எனவே
குறிச்சீட்டு~$l$ உடைய ஒரே கணுவை மட்டும் கொண்ட மரவுரு
$\tuple{0,l}$ ஆல் குறியீடாக்கப்படும்; குறிச்சீட்டு~$l_1$ உடைய
ஒரு வேருடன் இணைக்கப்பட்ட, குறிச்சீட்டுகள் $l_2$, $l_3$ உடைய இரு
தனிக் கணுக்களைக் கொண்ட மரவுரு $\tuple{2,
  \tuple{0, l_2}, \tuple{0, l_3}, l_1}$ ஆல் குறியீடாக்கப்படும்.

\begin{prop}
  \ollabel{prop:subtreeseq}
  குறியீடு~$t$ உடைய மரவுருவின் அனைத்து உட்மரவுருக்களின்
  குறியீடுகளை உறுப்புகளாகக் கொண்ட ஒரு தொடரின் குறியீட்டைத்
  திருப்பித்தரும் சார்பு $\fn{SubtreeSeq}(t)$ ஒரு முதன்மை
  மீள்வரையறைச் சார்பு.
\end{prop}

\begin{proof}
  $\fn{ISubtrees}(t) = \fn{subseq}(t, 1, (t)_0)$ ஒரு முதன்மை
  மீள்வரையறைச் சார்பு என்பதையும், மரவுரு~$t$ இன் உடனடி
  உட்மரவுருக்களின் குறியீடுகளை அது திருப்பித்தருகிறது என்பதையும்
  முதலில் கவனிக்க. இப்போது வேரிலிருந்து $n$ கணுக்கள் தொலைவிலுள்ள
  அனைத்து உட்மரவுருக்களின் தொடரைக் கணிக்கும் துணைச் சார்பு
  $\fn{hSubtreeSeq}(t,n)$ ஐ வரையறுக்கலாம். வேரிலிருந்து $0$
  கணுக்கள் தொலைவிலுள்ள~$t$ இன் உட்மரவுருக்களின் தொடர்---வேறுவிதமாகக்
  கூறினால்,~$t$ இன் வேரில் தொடங்குவது---~$t$ ஐ மட்டும் கொண்ட
  தொடராகும். $t$ இன் நிலை~$n+1$ இலுள்ள அனைத்து
  உட்மரவுருக்களின் தொடரைப் பெற, நிலை $n$ உட்மரவுருக்களுடன்,
  நிலை $n$ உட்மரவுருக்களின் அனைத்து உடனடி உட்மரவுருக்களையும் கொண்ட
  ஒரு தொடரைத் தொடரிணைக்கிறோம். இவை அனைத்தின் பட்டியலைப் பெற,
  $f(x)$ என்பது
  தொடர்களின் குறியீடுகளைத் திருப்பித்தரும் முதன்மை மீள்வரையறைச்
  சார்பு எனில், $g_f(s, k) = f((s)_0) \concat \dots \concat
  f((s)_k)$ உம் ஒரு முதன்மை மீள்வரையறைச் சார்பு என்பதைக் கவனிக்க:
    \begin{align*}
      g(s, 0) & = f((s)_0)\\
      g(s, k+1) & = g(s, k) \concat f((s)_{k+1})
    \end{align*}
    எடுத்துக்காட்டாக, $s$ ஒரு மரவுருக்களின் தொடர் எனில்,
    $h(s) = g_{\fn{ISubtrees}}(s, \len{s})$ ஆனது~$s$ இன்
    உறுப்புகளுடைய உடனடி உட்மரவுருக்களின் தொடரைத் தருகிறது. அதைப்
    பயன்படுத்தி $\fn{hSubtreeSeq}$ ஐப் பின்வருமாறு வரையறுக்கலாம்:
    \begin{align*}
      \fn{hSubtreeSeq}(t, 0) & = \tuple{t} \\
      \fn{hSubtreeSeq}(t, n+1) & = \fn{hSubtreeSeq}(t, n) \concat
      h(\fn{hSubtreeSeq}(t, n)).
    \end{align*}
    குறியீடு~$t$ உடைய மரவுருவிலுள்ள உட்மரவுருக்களின் மிகப்பெரிய
    நிலை, அதாவது வேருக்கும் ஓர் இலைக் கணுவுக்கும் இடையிலான
    மிகப்பெரிய தொலைவு, குறியீடு~$t$ ஆல் வரம்பிடப்படுகிறது.
    எனவே குறியீடு~$t$ உடைய மரவுருவின் அனைத்து உட்மரவுருக்களின்
    குறியீடுகளைக் கொண்ட ஒரு தொடர் $\fn{hSubtreeSeq}(t, t)$ ஆல்
    தரப்படுகிறது.
\end{proof}

\begin{prob}
  \olref[cmp][rec][tre]{prop:subtreeseq} இன் நிரூபிப்பிலுள்ள
  $\fn{hSubtreeSeq}$ இன் வரையறை பொதுவாக மீள்தோற்றங்களைக் கொண்டுள்ளது.
  விளைவுப் பட்டியலில் ஓர் உட்மரவுருவின் குறியீடு ஒருமுறை மட்டுமே
  தோன்றுவதை உறுதிசெய்யும் மாற்று வரையறையைத் தருக.
\end{prob}

\end{document}
